\section{Signed representatives and logarithmic general type}
\label{sec:signed-representatives}

Fix the terminal canonical counterexample $Y$ of
\Cref{prop:reduction}, with the very general locus supplied by
\Cref{prop:subvarieties}. The particular need for a signed divisor
comes from the slice bundle
$\mathbb P(\mathcal O_Y\oplus\mathcal O_Y(\iota K_Y))$.
Its two summands determine disjoint sections. Intersecting an integral
multisection of degree $e>0$, distinct from those sections, with them
and pushing to $Y$ gives effective divisors $D_0,D_\infty$ with
$D_0-D_\infty\sim e\iota K_Y$, as verified in
\Cref{eq:axis-linear-equivalence}. Thus the representative of $K_Y$
has signed coefficients, whereas the boundary is the reduced union
of the two supports, including any components that cancel in their
difference. We prove the logarithmic bigness needed to apply
\Cref{lem:degree-obstruction} to this boundary once its intersection
degree has been bounded.

The proposition allows any reduced boundary containing the signed
support. Its proof constructs an auxiliary nef adjoint with a
nonempty reduced floor. The analytic results in
\Cref{sec:analytic-inputs} lift a section from that whole floor;
only after this establishes nonnegative Kodaira dimension do we
apply the required semiampleness theorem.

\begin{proposition}[Signed representative alternative]
\label[proposition]{prop:signed-alternative}
Let $Y$ be a projective $\Q$-factorial terminal fourfold such that $K_Y$ is
nef and $\kappa(Y,K_Y)=-\infty$. Suppose that there is a very general locus
of $Y$ through whose points every proper positive-dimensional subvariety
is of general type on resolution. Let $J$ be a rational Weil divisor,
with coefficients of either sign, satisfying $J\sim_{\Q}K_Y$.
Let $p:W\to Y$ be a projective log resolution, and let $D_W$ be a reduced
simple normal crossings divisor containing the strict support of $J$
and every $p$-exceptional divisor. Then $K_W+D_W$ is big.
\end{proposition}

\begin{proof}
Put $P_W=K_W+D_W$. Terminality gives
\[
 K_W=p^*K_Y+E_W,\qquad E_W\geq0
\]
with $E_W$ exceptional. In particular, both $K_W$ and $P_W$ are
pseudo-effective. We assume that $P_W$ is not big and derive a
contradiction. We first reduce its Kodaira dimension to at most zero,
then construct a nef boundary with a nonempty reduced floor. We will
lift a nonzero section from that entire floor before using any
semiampleness theorem that requires nonnegative Kodaira dimension.

\subsection*{Positive Kodaira dimension would already imply bigness}

Suppose first that $\kappa(P_W)>0$. Resolve the rational map defined by
a moving subsystem of a sufficiently divisible multiple $bP_W$:
\[
 \pi:\widetilde W\longrightarrow W,\qquad
 f:\widetilde W\longrightarrow U.
\]
Here $\widetilde W$ is smooth, $U$ is the projective image, and, for a
very ample line bundle $\cO_U(1)$,
\begin{equation}
\label{eq:signed-moving-system}
 b\pi^*P_W\sim f^*\cO_U(1)+F,\qquad F\geq0.
\end{equation}
If $f$ is generically finite, then $P_W$ is big, contrary to the
assumption. Otherwise $0<\dim U<4$. A smooth fiber component through a
very general point of $\widetilde W$ is birational to a proper
positive-dimensional subvariety of $Y$ meeting the very general locus.
It is therefore of general type. We work over a smooth open of $U$
where generic smoothness applies. On its smooth fibers the restriction
of $K_{\widetilde W}$ is the canonical bundle of the fiber, up to the
constant one-dimensional factor coming from the base.

Fix an ample Cartier divisor $H_0$ on $\widetilde W$. For some integer
$j>0$, the restriction of $jK_{\widetilde W}-H_0$ has a nonzero section
on such a fiber component. Distinct components of a smooth fiber are
disjoint, so the section can be taken to be zero on the other
components. Choose the fiber over a point where the fiber dimension of
the space of sections is the generic one for every integer $j$; this
requires only countably many exclusions. Consequently
\[
 f_*\cO_{\widetilde W}(jK_{\widetilde W}-H_0)
\]
is generically nonzero for one $j$. After tensoring by a sufficiently
large power $\cO_U(c)$, this coherent sheaf has a nonzero global
section. Thus
\[
 jK_{\widetilde W}+c f^*\cO_U(1)-H_0\sim E_0,
 \qquad E_0\geq0.
\]
The divisor on the left before subtracting $H_0$ is big. By
\Cref{eq:signed-moving-system}, adding an effective divisor gives
bigness of $jK_{\widetilde W}+cb\pi^*P_W$. Birational pushforward gives
bigness of $jK_W+cbP_W$. For completeness, the ample divisor $H_0$
has big pushforward: for a fixed ample divisor $H_W$ on $W$ and small
positive rational $\eta$, $H_0-\eta\pi^*H_W$ is ample, and pushing an
effective rational representative proves this assertion. Finally
\[
 (j+cb)P_W=jK_W+cbP_W+jD_W
\]
is big. This contradiction proves
\begin{equation}
\label{eq:signed-kodaira-small}
 \kappa(P_W)\leq0.
\end{equation}

\subsection*{A minimal model and a crepant klt perturbation}

Run a log minimal model program for the rational dlt pair $(W,D_W)$.
We use flip existence for klt pairs, the usual dlt contraction
properties, and termination of flips for pseudo-effective lc
fourfolds, applied with zero nef part
\cite{BCHM2010,ChenTsakanikas2023}. The dlt flips needed here are
obtained by lowering the finitely many boundary coefficients slightly
while retaining negativity of the given extremal ray. The resulting
pair is klt. On the relative Picard-number-one contraction, the
original and perturbed negative adjoints are relatively proportional
with positive rational ratio; the contraction theorem descends a
Cartier multiple of their numerically trivial difference. The klt
flip therefore has the required positivity for the original adjoint
as well. This supplies the flips used in this dlt program.

There are only finitely many divisorial contractions. Pseudo-effectivity
persists under each pushforward and strict transform and excludes a
negative fiber space, by intersection with covering curves in a
general fiber. These pushforwards on numerical divisor classes are
well-defined on the $\Q$-factorial models: on a common resolution they
are pullback followed by pushforward, and exceptional divisor classes
give no ambiguity. We obtain a projective $\Q$-factorial dlt minimal
model $(Z,D_Z)$ without extracting divisors. Write
\[
 M=K_Z+D_Z.
\]
The divisor $D_Z$ is reduced, $M$ is nef, and the standard exceptional
comparison on a common resolution preserves section spaces in
sufficiently divisible degrees. Hence
\[
 \kappa(M)=\kappa(P_W)\leq0.
\]
Moreover $K_Z$ is pseudo-effective, since $K_W$ is pseudo-effective.

The signed representative supplies an additional linear relation.
Indeed
\[
 P_W\sim_{\Q}p^*J+E_W+D_W,
\]
and the divisor on the right is supported on $D_W$. Its strict
pushforward to $Z$ is a signed rational divisor $G_Z$ supported on
$D_Z$, with
\begin{equation}
\label{eq:signed-supported-representative}
 M\sim_{\Q}G_Z.
\end{equation}

We next make a small klt perturbation while keeping a fixed Cartier
multiple of $M$. Choose an integer $\ell>0$ for which $\ell M$ is
Cartier, and fix a rational number
\begin{equation}
\label{eq:signed-delta-choice}
 0<\delta<\min\left\{\frac12,\frac{1}{16\ell+2}\right\}.
\end{equation}
The klt adjoint
\[
 M-\delta D_Z=K_Z+(1-\delta)D_Z
\]
is pseudo-effective because $K_Z$ is pseudo-effective and
$(1-\delta)D_Z$ is effective. Run its minimal model program. We show
inductively that every step is crepant for $M$, preserves nefness of
$M$, and preserves the Cartier property of $\ell M$.

At a given step let $R$ be a negative extremal ray for the driving
adjoint. Since $M$ is nef and $\delta<1/2$, the ray is also negative
for the klt adjoint $M-D_Z/2$. The length bound gives a rational curve
$C$ spanning $R$ such that
\[
 0<-(M-D_Z/2)\cdot C\leq 2\dim Z=8.
\]
Set $a_C=M\cdot C$ and $b_C=D_Z\cdot C$. Then
\[
 a_C\geq0,\qquad b_C\leq2a_C+16,
 \qquad a_C<\delta b_C.
\]
It follows that
\[
 0\leq a_C<\frac{16\delta}{1-2\delta}<\frac1\ell.
\]
Because $\ell M$ is Cartier, $a_C\in\ell^{-1}\Z$, and therefore
$a_C=0$.

The line-bundle clause of the contraction theorem now descends
$\cO_Z(\ell M)$ to an actual line bundle on the contraction base
\cite[Theorem~1.1(4)(iii)]{Fujino2011}; the same theorem's part~(5)
gives the length bound just used. For a divisorial contraction its
pullback is the original bundle. For a flip it pulls back on both
sides. Consequently $M$ remains nef, the same multiple $\ell M$ remains Cartier, and its pullbacks to a common resolution agree. This is
crepancy for the full adjoint, rather than numerical triviality
alone. More explicitly, choose a rational section of the descended
line bundle. The divisor $\ell M$ differs from its pullback divisor
by a principal divisor; transform the same rational function across
the flip. The two exact divisor pullbacks then agree on a common
resolution, with compatible canonical divisors. The driving pair
stays klt. The boundary with coefficient
$1/2$ is smaller than the driving boundary, so it is again klt at
the next step. The preceding argument therefore repeats with the
same $\ell$ and $\delta$.

Let $Z'$ be the resulting model, and denote strict transforms by
$D'$, $M'$, and $G'$. Put $K'=K_{Z'}$. We have
\begin{equation}
\label{eq:signed-nef-interval}
 \begin{gathered}
 (Z',D')\text{ is lc},\qquad K'\text{ is pseudo-effective},
 \qquad M'\sim_{\Q}G',\qquad \kappa(M')\leq0,\\
 M'-uD'\text{ is nef and }(Z',(1-u)D')\text{ is klt}
 \quad(0<u\leq\delta).
 \end{gathered}
\end{equation}
Here the last nefness assertion follows by interpolation between
$M'$ and $M'-\delta D'$. The klt assertion follows by interpolation
between the klt driving pair and the lc full-boundary pair. The
underlying variety $Z'$ is klt.

We have reduced the problem to a fixed nef interval of actual klt
adjoints and a signed representative supported on its boundary.
We next use that representative to locate a nonempty boundary floor.

\subsection*{A nonempty floor and the reduced non-klt ideal}

Some coefficient of $G'$ is positive. Suppose otherwise. Intersect
$K'+D'\sim_{\Q}G'\leq0$ with the third power of an ample divisor.
Pseudo-effectivity of $K'$ and effectivity of $D'$ show that all
inequalities are equalities; positivity for nonzero effective
divisors then gives $D'=0$ and $G'=0$.

To see why this is impossible, take a common smooth resolution $V$
of $W$ and $Z'$. The pullback to $V$ of the signed divisor $J$ is
supported over $D_W$. Since the programs extract no divisors and
$D'=0$, this pullback is exceptional over $Z'$. It is rationally
linearly equivalent to the pullback of $K_Y$, and hence is nef.
The negativity lemma makes this exceptional divisor nonpositive.
Its intersection with an ample divisor to the third power is
nonnegative by nefness and nonpositive by its sign; it must vanish.
Thus the pullback of $K_Y$ is rationally linearly trivial, contrary
to $\kappa(Y,K_Y)=-\infty$.

Write $D'=\sum_j D'_j$ and $G'=\sum_j g_jD'_j$, allowing $g_j=0$.
Let $a=\max_j g_j>0$. Choose a sufficiently small positive rational
$s$ and define
\begin{equation}
\label{eq:signed-boundary-wall}
 B=(1-s)D'+\frac{s}{a}G',\qquad
 \alpha=1+\frac{s}{a},\qquad N=K'+B.
\end{equation}
Choose $s$ so that every coefficient $1-s+(s/a)g_j$ is nonnegative
and $s/\alpha\in(0,\delta)$. All these coefficients are at most one,
with equality precisely when $g_j=a$. Therefore
\[
 0\leq B\leq D',\qquad S:=\floor{B}\ne0,
 \qquad N\sim_{\Q}\alpha M'-sD'.
\]
The divisor $N$ is nef by \Cref{eq:signed-nef-interval}. The pair
$(Z',B)$ is lc and klt outside $S$. Indeed, outside $S$ its finitely
many boundary coefficients are bounded above by a number less than
one; interpolate the lc full-boundary pair with the underlying klt
variety. The same reasoning shows that $(Z',bB)$ is klt whenever
$0<b<1$.

The multiplier ideals are actual ideals
\begin{equation}
\label{eq:signed-multiplier-ideals}
 \cJ(Z',B)=\cI_S,\qquad
 \cJ(Z',bB)=\cO_{Z'}\quad(0<b<1).
\end{equation}
Here $S$ carries its reduced induced structure. To verify the first
identity, use the discrepancy description on a log resolution.
Since the pair is lc, a regular function satisfies every positive
log-discrepancy condition automatically. At a log-discrepancy-zero
divisor it must vanish to positive order. All such centers lie in
$S$, and each component of $S$ itself gives such a condition. A
function vanishing on reduced $S$ has positive order at every
valuation centered in $S$; conversely the component valuations
force vanishing on $S$. This proves the identity. These are
multiplier ideals for an lc pair, not a substitution of a non-lc
ideal for the reduced ideal.

\subsection*{Lifting sections from the floor}

We claim that, for every sufficiently large and divisible integer $m$,
the following restriction map to the whole reduced floor $S$ is surjective:
\begin{equation}
\label{eq:signed-section-surjection}
 H^0(Z',\cO_{Z'}(mN))\longrightarrow H^0(S,\cO_S(mN)).
\end{equation}
The interior injectivity theorem will compare the boundary $B$
with $b_0B$ for $b_0<1$, where the pair is klt. It also requires
an endpoint $b_2B$ beyond $B$. The corresponding residual bundle
is $(m-1)N-(b_2-1)B$; its positivity is supplied by the interval
\Cref{eq:signed-nef-interval}. We now verify this for both endpoints.

Fix rational numbers $0<b_0<1=b_1<b_2$ and a projective
log resolution $h:R\to Z'$ of $(Z',D')$, hence also of $(Z',B)$.
Choose a fixed
effective integral $h$-exceptional divisor $A_R$, with coefficients
large enough that
\[
 C_i=A_R+h^*(K'+b_iB)-K_R\geq0\quad(i=0,1,2).
\]
Their supports lie in one simple normal crossings divisor, and
$C_0\leq C_1\leq C_2$. Moreover $C_1$ is the strict interior convex
combination of $C_0$ and $C_2$ with parameter
$(1-b_0)/(b_2-b_0)$. For $mN$ Cartier put
\[
 L_R=A_R+h^*(mN)-K_R.
\]
This is an integral divisor. Its endpoint residuals are
\begin{align}
 L_R-C_i
 &=h^*((m-1)N+(1-b_i)B) \notag\\
 &\sim_{\Q}h^*\bigl(t_i(m)M'-v_i(m)D'\bigr),
 \label{eq:signed-endpoint-residuals}\\
 t_i(m)&=(m-1)\alpha+(1-b_i)s/a,\qquad
 v_i(m)=(m-1)s-(1-b_i)(1-s). \notag
\end{align}
For $i=0,2$, $t_i(m)>0$ once $m$ is large, and
\[
 u_i(m):=\frac{v_i(m)}{t_i(m)}\longrightarrow\frac{s}{\alpha}
 \in(0,\delta).
\]
Consequently each endpoint residual is rationally linearly
equivalent to a positive rational multiple of
$h^*(K'+(1-u_i(m))D')$, a nef klt adjoint pulled back to its log
resolution. By \Cref{thm:adjoint-metric}, it admits a semipositive
singular metric with zero Lelong numbers at every point. Taking
positive rational powers and transporting by the actual rational
line-bundle isomorphisms preserves this property. The endpoint
hypotheses of \Cref{thm:interior-injectivity} are therefore satisfied.

That theorem gives injectivity from the first ordinary cohomology
group of $K_R+L_R-\floor{C_1}$ to that of
$K_R+L_R-\floor{C_0}$. The precise floor calculation is
\begin{equation}
\label{eq:signed-floor-calculation}
 K_R+L_R-\floor{C_i}
 =h^*(mN)+\ceil{K_R-h^*(K'+b_iB)}.
\end{equation}
Multiplier-ideal local vanishing and the projection formula
\cite[Theorem~3.3]{Fujino2011} identify these groups with
\[
 H^1\bigl(Z',\cO_{Z'}(mN)\otimes\cJ(Z',b_iB)\bigr).
\]
The comparison map is the map induced by inclusion of the
multiplier ideals. By \Cref{eq:signed-multiplier-ideals}, we have
proved injectivity of
\[
 H^1(Z',\cO_{Z'}(mN)\otimes\cI_S)
 \longrightarrow H^1(Z',\cO_{Z'}(mN)).
\]
The exact sequence of $S$ proves
\Cref{eq:signed-section-surjection}.

The restriction map is now surjective. To obtain a nonzero global
section, we must produce one on the whole reduced scheme $S$.
We do this on a dlt floor using adjunction and threefold abundance.
The next step therefore includes both gluing on that floor and
actual line-bundle descent to $S$.

\subsection*{Whole-floor abundance and descent}

Take a projective $\Q$-factorial dlt blowup
\[
 g:(Q,B_Q)\longrightarrow(Z',B),\qquad
 K_Q+B_Q=g^*(K'+B)=g^*N,
\]
extracting only lc places; this is the lc case of the dlt blowup
theorem \cite[Theorem~2.10]{FujinoHashizume2023}. Set
$T=\floor{B_Q}$. It is nonempty. Whole-floor dlt adjunction equips
$T$ with an effective rational different $\operatorname{Diff}_T$
such that
\begin{equation}
\label{eq:signed-whole-floor-adjunction}
 K_T+\operatorname{Diff}_T
 =(K_Q+B_Q)|_T=g^*N|_T
\end{equation}
as the actual rational adjoint line bundle, including the conductor
and residue identifications on the normalization. The pair on $T$
is semi-dlt; see \cite[Remark~1.2(3)]{Fujino2000} and
\cite[Remark~2.7]{FujinoGongyo2014}.

We record the depth condition needed for this whole-floor statement.
The underlying $Q$ is klt and $\Q$-factorial. Both $\cO_Q$ and
$\cO_Q(-T)$ are Cohen--Macaulay. For the latter assertion, locally
trivialize a Cartier multiple of $T$ and take the normal cyclic index
cover. It is finite and quasi-\'etale, hence klt and Cohen--Macaulay.
Its finite pushforward decomposes into the corresponding rank-one
reflexive sheaves, including $\cO_Q(-T)$; in characteristic zero
these are direct summands. They are therefore Cohen--Macaulay.
The sequence
\[
 0\longrightarrow\cO_Q(-T)\longrightarrow\cO_Q
 \longrightarrow\cO_T\longrightarrow0
\]
then gives $S_2$ for $T$ (in fact the required Cohen--Macaulay depth
along its support). Its codimension-one crossings and the conductor
description come from the dlt structure. Thus the semi-log-canonical
adjunction in \Cref{eq:signed-whole-floor-adjunction} is on the
whole projective threefold, not on a disjoint collection of its
components.

The adjoint in \Cref{eq:signed-whole-floor-adjunction} is nef.
Projective semi-dlt threefold abundance makes it semiample
\cite[Corollary~4.10]{Fujino2000}. Hence
\[
 H^0(T,\cO_T(mg^*N))\ne0
\]
for all sufficiently divisible positive $m$.

These sections descend to $S$. Indeed
$\cJ(Q,B_Q)=\cO_Q(-T)$, and crepancy together with multiplier-ideal
local vanishing gives
\begin{equation}
\label{eq:signed-floor-pushforward}
 g_*\cO_Q(-T)=\cI_S,\qquad
 R^i g_*\cO_Q(-T)=0\quad(i>0).
\end{equation}
One can compute both assertions on a common log resolution of
$(Q,B_Q)$ and $(Z',B)$. The equality of their pulled-back adjoints
identifies the discrepancy round-up sheaves, and local vanishing
for each log resolution followed by Leray gives
\Cref{eq:signed-floor-pushforward}. Since $Z'$ is normal,
$g_*\cO_Q=\cO_{Z'}$. Pushing forward the divisor sequence therefore
gives
\[
 g_*\cO_T=\cO_S.
\]
The projection formula identifies the nonzero section spaces above
with $H^0(S,\cO_S(mN))$. Choose $m$ both sufficiently divisible
for this semiampleness and sufficiently large for
\Cref{eq:signed-section-surjection}. A nonzero section on $T$
then descends to $S$ and lifts to $Z'$. We have proved
$\kappa(N)\geq0$.

\subsection*{The final klt adjoint}

Because $B\leq D'$, section inclusion in divisible degrees gives
$\kappa(N)\leq\kappa(M')\leq0$. Thus $\kappa(N)=0$.
Now consider the klt pair
\[
 \left(Z',\left(1-\frac{s}{\alpha}\right)D'\right).
\]
Its nef adjoint is rationally linearly equivalent to $N/\alpha$,
so its Kodaira dimension is zero. By
\Cref{thm:adjoint-metric}, on a smooth projective log resolution
a positive Cartier multiple of its pullback has a semipositive
singular metric with zero point Lelong numbers. All hypotheses of
the fourfold semiampleness result of Gongyo--Matsumura are now
satisfied: the pair is projective klt with rational boundary, the
adjoint has nonnegative Kodaira dimension, and the specified
pullback metric exists
\cite[Corollary~5.3]{GongyoMatsumura2017}.

Semiampleness and Kodaira dimension zero make this adjoint
numerically trivial. On the other hand, $K'$ is pseudo-effective,
$D'\ne0$ because $G'$ has a positive coefficient, and
$1-s/\alpha>0$. Intersecting
$K'+(1-s/\alpha)D'$ with the third power of an ample divisor gives
a strictly positive number. This contradiction proves the
proposition.
\end{proof}
